\section{习题}

\subsection{本章速览}

\begin{itemize}
  \item \textbf{基本概念：}方程的阶是最高阶导数的阶数；未知函数及其各阶导数都是一次的叫线性方程。$n$ 阶方程的通解含 $n$ 个独立任意常数，初值条件也要给 $n$ 个。
  \item \textbf{可分离方程：}能写成 $y'=f(x)g(y)$ 的方程，分离后两边积分。除以 $g(y)$ 时可能丢掉 $g(y)=0$ 对应的常值解，最后要单独检查。
  \item \textbf{一阶线性方程：}先化成标准形 $y'+p(x)y=q(x)$，取积分因子 $\mu(x)=e^{\int p(x)\mathrm{d}x}$，则 $(\mu y)'=\mu q$，通解为 $y=\mu^{-1}\bigl(\int \mu q\,\mathrm{d}x+C\bigr)$。要求 $p,q$ 在讨论区间上连续。
  \item \textbf{齐次方程：}形如 $y'=\varphi(y/x)$，令 $u=y/x$，则 $y=ux$、$y'=u+xu'$，方程化为可分离的 $xu'=\varphi(u)-u$。
  \item \textbf{常系数线性方程：}齐次方程由特征方程 $r^2+pr+q=0$ 的根决定通解（两不等实根、重根、共轭复根三种情形）。非齐次方程右端为 $e^{\lambda x}P_n(x)$ 时，若 $\lambda$ 是 $k$ 重特征根，特解应设为 $x^kQ_n(x)e^{\lambda x}$，其中 $Q_n$ 为 $n$ 次待定多项式。
  \item \textbf{欧拉方程与积分方程：}欧拉方程 $x^2y''+axy'+by=f(x)$ 在 $x>0$ 时令 $x=e^t$，有 $xy'=Dy$、$x^2y''=D(D-1)y$（$D=\mathrm{d}/\mathrm{d}t$），化为常系数方程。含变限积分的方程先拆开积分号再求导，化为微分方程；每求导一次，就要由原方程补充一个初始条件。
\end{itemize}

\subsection{基础过关题}

\textbf{1.} 求解初值问题
\[
\frac{\mathrm{d}y}{\mathrm{d}x}=y(y-1),\qquad y(0)=2,
\]
并指出解的最大存在区间。

\ifshowsolutions%
\boxed{\textbf{解}}
\textbf{思路：}右端只含 $y$，属于可分离方程。先分离变量，再代入初值定常数；$y\equiv 0$ 与 $y\equiv 1$ 是除法时可能丢掉的常值解，需单独检查。

当 $y\neq 0,1$ 时，
\[
\frac{\mathrm{d}y}{y(y-1)}=\mathrm{d}x.
\]
而
\[
\frac{1}{y(y-1)}=\frac{1}{y-1}-\frac{1}{y},
\]
两边积分得
\[
\ln\left|\frac{y-1}{y}\right|=x+C.
\]
代入 $x=0,\ y=2$：
\[
\ln\frac12=C.
\]
由于 $y(0)=2>1$，在含 $0$ 的区间上 $y>1$，故可去掉绝对值：
\[
\frac{y-1}{y}=\frac{1}{2}e^{x}.
\]
解出
\[
y\left(1-\frac{e^x}{2}\right)=1
\quad\Longrightarrow\quad
y=\frac{2}{2-e^x}.
\]
常值解 $y\equiv 0,1$ 不满足 $y(0)=2$，故不影响本题。解要有意义需 $2-e^x>0$，即 $x<\ln 2$；当 $x\to(\ln 2)^-$ 时 $y\to+\infty$。

\textbf{答案：}\boxed{y=\frac{2}{2-e^x},\quad x\in(-\infty,\ln 2).}

\else \answerarea[8] \fi%
\textbf{2.} 求初值问题
\[
y'+y\tan x=\sec x,\qquad y(0)=0,\qquad -\frac{\pi}{2}<x<\frac{\pi}{2}
\]
的解。

\ifshowsolutions%
\boxed{\textbf{解}}
\textbf{思路：}这是一阶线性方程，$p(x)=\tan x$，$q(x)=\sec x$。在给定区间上 $\cos x>0$，可直接取积分因子 $\mu=e^{\int \tan x\,\mathrm{d}x}$。

因为
\[
\int \tan x\,\mathrm{d}x=-\ln(\cos x),
\]
所以
\[
\mu(x)=e^{-\ln(\cos x)}=\sec x.
\]
方程两边乘 $\sec x$：
\[
y'\sec x+y\sec x\tan x=\sec^2 x.
\]
注意 $(\sec x)'=\sec x\tan x$，故左边为 $(y\sec x)'$，于是
\[
(y\sec x)'=\sec^2 x.
\]
积分得
\[
y\sec x=\tan x+C,
\]
即
\[
y=\sin x+C\cos x.
\]
代入 $y(0)=0$，得 $C=0$。

代回检验：$y'=\cos x$，
\[
y'+y\tan x=\cos x+\sin x\tan x
=\frac{\cos^2x+\sin^2x}{\cos x}=\sec x.
\]

\textbf{答案：}\boxed{y=\sin x,\quad -\frac{\pi}{2}<x<\frac{\pi}{2}.}

\else \answerarea[8] \fi%
\textbf{3.} 求解初值问题
\[
y'=\frac{x^2+y^2}{2xy},\qquad y(2)=1,\qquad x>0.
\]

\ifshowsolutions%
\boxed{\textbf{解}}
\textbf{思路：}右端分子、分母都是二次齐次式，可写成 $y/x$ 的函数，因此令 $u=y/x$。

令 $u=\dfrac{y}{x}$，则 $y=ux$，$y'=u+xu'$。原方程化为
\[
u+xu'=\frac{1+u^2}{2u}.
\]
移项：
\[
xu'=\frac{1+u^2}{2u}-u=\frac{1-u^2}{2u}.
\]
当 $u\neq\pm1$ 时分离变量：
\[
\frac{2u}{1-u^2}\,\mathrm{d}u=\frac{\mathrm{d}x}{x}.
\]
积分：
\[
-\ln|1-u^2|=\ln x+C_1,
\]
即
\[
x(1-u^2)=K.
\]
代入 $x=2,\ u=y/x=1/2$：
\[
K=2\left(1-\frac14\right)=\frac32.
\]
所以
\[
x\left(1-\frac{y^2}{x^2}\right)=\frac32
\quad\Longrightarrow\quad
x^2-y^2=\frac32x.
\]
由 $y(2)=1>0$，取正分支：
\[
y=\sqrt{x^2-\frac32x}.
\]
再检查 $u=\pm1$ 即 $y=\pm x$：在 $x=2$ 处分别取 $\pm2$，不满足 $y(2)=1$，故不是本题的解。根号内需 $x^2-\frac32x>0$，结合 $x>0$ 得 $x>\frac32$。

\textbf{答案：}\boxed{y=\sqrt{x^2-\frac32x},\quad x\in\left(\frac32,+\infty\right).}

\else \answerarea[8] \fi%
\textbf{4.} 求初值问题
\[
y''-3y'+2y=0,\qquad y(0)=1,\qquad y'(0)=0
\]
的解。

\ifshowsolutions%
\boxed{\textbf{解}}
\textbf{思路：}二阶常系数齐次线性方程，先写特征方程，再按根的情形写通解，最后用两个初值条件定常数。

特征方程为
\[
r^2-3r+2=0
\quad\Longrightarrow\quad
(r-1)(r-2)=0,
\]
故 $r_1=1,\ r_2=2$，通解为
\[
y=C_1e^x+C_2e^{2x}.
\]
由 $y(0)=1$ 得
\[
C_1+C_2=1.
\]
又
\[
y'=C_1e^x+2C_2e^{2x},
\]
由 $y'(0)=0$ 得
\[
C_1+2C_2=0.
\]
两式相减得 $C_2=-1$，从而 $C_1=2$。

\textbf{答案：}\boxed{y=2e^x-e^{2x}.}

\else \answerarea[8] \fi%
\textbf{5.} 求欧拉方程
\[
x^2y''-3xy'+4y=0\qquad (x>0)
\]
的通解。

\ifshowsolutions%
\boxed{\textbf{解}}
\textbf{思路：}欧拉方程可令 $y=x^r$ 直接代入，也可令 $x=e^t$ 化为常系数方程。本题用 $y=x^r$ 更快。

设 $y=x^r$，则
\[
y'=rx^{r-1},\qquad y''=r(r-1)x^{r-2}.
\]
代入原方程并约去 $x^r$：
\[
r(r-1)-3r+4=0,
\]
即
\[
r^2-4r+4=(r-2)^2=0.
\]
特征根 $r=2$ 为二重根，因此
\[
y=(C_1+C_2\ln x)x^2.
\]
也可用 $x=e^t$ 验证：$D(D-1)y-3Dy+4y=(D-2)^2y=0$，通解为 $e^{2t}(C_1+C_2t)$，代回 $t=\ln x$ 即得上式。

\textbf{答案：}\boxed{y=x^2(C_1+C_2\ln x),\quad x>0.}

\else \answerarea[8] \fi%
\subsection{综合提高题}

\textbf{6.} 求初值问题
\[
y''-4y'+4y=e^{2x},\qquad y(0)=1,\qquad y'(0)=0
\]
的解。

\ifshowsolutions%
\boxed{\textbf{解}}
\textbf{思路：}齐次部分的特征根为二重根 $2$，右端 $e^{2x}$ 中的指数 $2$ 恰好也是二重特征根，属于“共振”情形，特解要乘 $x^2$。

齐次方程 $y''-4y'+4y=0$ 的特征方程为
\[
r^2-4r+4=(r-2)^2=0,
\]
故
\[
y_h=(C_1+C_2x)e^{2x}.
\]
设特解 $y^*=Ax^2e^{2x}$。利用
\[
(D-2)^2\bigl[e^{2x}v(x)\bigr]=e^{2x}v''(x),
\]
取 $v(x)=Ax^2$，得
\[
(D-2)^2(Ax^2e^{2x})=e^{2x}(Ax^2)''=2Ae^{2x}.
\]
令其等于 $e^{2x}$，得 $A=\frac12$。于是
\[
y=e^{2x}\left(C_1+C_2x+\frac{x^2}{2}\right).
\]
由 $y(0)=1$ 得 $C_1=1$。求导：
\[
y'=e^{2x}\left[2\left(C_1+C_2x+\frac{x^2}{2}\right)+\left(C_2+x\right)\right]
=e^{2x}\left[2C_1+C_2+(2C_2+1)x+x^2\right].
\]
由 $y'(0)=0$ 得
\[
2C_1+C_2=0
\quad\Longrightarrow\quad
C_2=-2.
\]

\textbf{答案：}\boxed{y=e^{2x}\left(1-2x+\frac{x^2}{2}\right).}

\else \answerarea[10] \fi%
\textbf{7.} 求欧拉方程
\[
x^2y''+xy'-y=\ln x\qquad (x>0)
\]
的通解。

\ifshowsolutions%
\boxed{\textbf{解}}
\textbf{思路：}先作代换 $x=e^t$ 把欧拉方程化为常系数方程，再求齐次通解和一个特解。

令 $x=e^t$，即 $t=\ln x$，记 $D=\dfrac{\mathrm{d}}{\mathrm{d}t}$。由
\[
xy'=Dy,\qquad x^2y''=D(D-1)y,
\]
原方程化为
\[
D(D-1)y+Dy-y=t,
\]
即
\[
(D^2-1)y=t.
\]
对应齐次方程 $D^2y-y=0$ 的特征方程为 $r^2-1=0$，根为 $r=\pm1$，故
\[
y_h=C_1e^t+C_2e^{-t}=C_1x+\frac{C_2}{x}.
\]
对右端 $t$，注意到 $D^2(-t)=0$，而 $-( -t)=t$，所以 $y_p=-t$ 是一个特解。因此
\[
y=C_1e^t+C_2e^{-t}-t.
\]
代回 $t=\ln x$ 即得通解。直接验证 $y=-\ln x$：
\[
x^2\left(\frac{1}{x^2}\right)+x\left(-\frac1x\right)-(-\ln x)
=1-1+\ln x=\ln x.
\]

\textbf{答案：}\boxed{y=C_1x+\frac{C_2}{x}-\ln x,\quad x>0.}

\else \answerarea[10] \fi%
\textbf{8.} 设 $f(x)$ 在 $[0,+\infty)$ 上连续，且满足积分方程
\[
f(x)=e^x+\int_0^x (x-t)f(t)\,\mathrm{d}t,
\]
求 $f(x)$。

\ifshowsolutions%
\boxed{\textbf{解}}
\textbf{思路：}先把积分号内的 $(x-t)$ 拆开，使 $x$ 移到积分号外，再对变上限积分求导，把积分方程化为二阶常系数线性方程。初始条件由原方程及一阶导方程在 $x=0$ 处给出。

将积分拆开：
\[
f(x)=e^x+x\int_0^x f(t)\,\mathrm{d}t-\int_0^x tf(t)\,\mathrm{d}t.
\]
因 $f$ 连续，右边可导，求导得
\[
f'(x)=e^x+\int_0^x f(t)\,\mathrm{d}t+xf(x)-xf(x)
=e^x+\int_0^x f(t)\,\mathrm{d}t.
\]
再求导：
\[
f''(x)=e^x+f(x),
\]
即
\[
f''-f=e^x.
\]
由原方程令 $x=0$：
\[
f(0)=e^0+0=1.
\]
由 $f'(x)=e^x+\int_0^x f(t)\,\mathrm{d}t$ 令 $x=0$：
\[
f'(0)=1.
\]
齐次方程 $f''-f=0$ 的通解为 $C_1e^x+C_2e^{-x}$。右端 $e^x$ 对应的指数 $1$ 是单特征根，设特解 $Axe^x$。因为
\[
(D^2-1)(Axe^x)=2Ae^x,
\]
故 $A=\frac12$，从而
\[
f(x)=\left(C_1+\frac{x}{2}\right)e^x+C_2e^{-x}.
\]
由 $f(0)=1$ 得
\[
C_1+C_2=1.
\]
求导：
\[
f'(x)=\left(C_1+\frac{x}{2}+\frac12\right)e^x-C_2e^{-x}.
\]
由 $f'(0)=1$ 得
\[
C_1+\frac12-C_2=1
\quad\Longrightarrow\quad
C_1-C_2=\frac12.
\]
解得 $C_1=\frac34,\ C_2=\frac14$。

\textbf{答案：}\boxed{f(x)=\left(\frac34+\frac{x}{2}\right)e^x+\frac14e^{-x}.}

\else \answerarea[10] \fi%
\textbf{9.} 设曲线 $y=y(x)\ (x>0)$ 过点 $(1,1)$，且该曲线上任一点 $(x,y)$ 处的切线在 $y$ 轴上的截距等于 $x^2$，求该曲线的方程。

\ifshowsolutions%
\boxed{\textbf{解}}
\textbf{思路：}先写出切线方程，求出它在 $y$ 轴上的截距，由题意列出微分方程，再按一阶线性方程求解，最后用点 $(1,1)$ 定常数。

曲线在点 $(x,y)$ 处的切线为
\[
Y-y=y'(X-x).
\]
令 $X=0$，得切线在 $y$ 轴上的截距
\[
Y=y-xy'.
\]
由题意
\[
y-xy'=x^2.
\]
两边除以 $-x$（$x>0$），化成标准形：
\[
y'-\frac1xy=-x.
\]
取积分因子
\[
\mu(x)=e^{\int(-1/x)\mathrm{d}x}=e^{-\ln x}=\frac1x.
\]
方程两边乘 $\frac1x$：
\[
\frac{y'}{x}-\frac{y}{x^2}=-1
\quad\Longrightarrow\quad
\left(\frac{y}{x}\right)'=-1.
\]
积分得
\[
\frac{y}{x}=-x+C,
\]
即
\[
y=Cx-x^2.
\]
代入 $y(1)=1$：
\[
C-1=1\quad\Longrightarrow\quad C=2.
\]
检验：$y=2x-x^2$，$y'=2-2x$，切线在 $y$ 轴上的截距为
\[
y-xy'=(2x-x^2)-x(2-2x)=x^2,
\]
符合题意。

\textbf{答案：}\boxed{y=2x-x^2,\quad x>0.}

\else \answerarea[8] \fi%
\subsection{真题风格与综合训练}

\textbf{10.}（选择题）设 $p(x),q(x)$ 在区间 $I$ 上连续，$q(x)\not\equiv0$，$y_1(x),y_2(x)$ 是方程
\[
y'+p(x)y=q(x)
\]
在 $I$ 上的两个不同的解，则下列函数一定是方程 $y'+p(x)y=0$ 的解的是（\quad）。

\begin{center}
\begin{tabular}{ll}
(A) $y_1+y_2$ & (B) $y_1-y_2$ \\[4pt]
(C) $2y_1-y_2$ & (D) $y_1y_2$
\end{tabular}
\end{center}

\ifshowsolutions%
\boxed{\textbf{解}}
\textbf{思路：}记 $L[y]=y'+p(x)y$，则 $L$ 是线性算子，且 $L[y_1]=L[y_2]=q(x)$。

对四个选项分别计算：
\[
L[y_1-y_2]=L[y_1]-L[y_2]=q-q=0,
\]
所以 $y_1-y_2$ 是齐次方程的解。
其余选项：
\[
L[y_1+y_2]=2q\neq 0,\qquad
L[2y_1-y_2]=2q-q=q\neq 0
\]
（这里用到 $q(x)\not\equiv0$），而 $L[y_1y_2]$ 一般不等于 $0$，不能保证是齐次方程的解。

\textbf{答案：}\boxed{(B)}

\else \answerarea[4] \fi%
\textbf{11.}（填空题）已知二阶常系数齐次线性方程
\[
y''+ay'+by=0
\]
的通解为
\[
y=e^{-x}\left(C_1\cos 2x+C_2\sin 2x\right),
\]
则 $a=\underline{\qquad}$，$b=\underline{\qquad}$。

\ifshowsolutions%
\boxed{\textbf{解}}
\textbf{思路：}由通解看出特征根为 $-1\pm 2i$，写出对应的特征方程，再与 $r^2+ar+b=0$ 比较系数。

通解中的指数因子 $e^{-x}$ 给出实部 $-1$，三角函数 $\cos 2x,\sin 2x$ 给出虚部 $\pm2$，故特征根为
\[
r_{1,2}=-1\pm 2i.
\]
特征方程为
\[
(r+1)^2+4=0,
\]
即
\[
r^2+2r+5=0.
\]
与 $r^2+ar+b=0$ 比较，得 $a=2,\ b=5$。

\textbf{答案：}\boxed{a=2,\quad b=5.}

\else \answerarea[4] \fi%
\textbf{12.}（解答题）设函数 $y=y(x)$ 满足
\[
y''+2y'+5y=0,\qquad y(0)=1,\qquad y'(0)=0.
\]
求广义积分
\[
I=\int_0^{+\infty}y(x)\,\mathrm{d}x.
\]

\ifshowsolutions%
\boxed{\textbf{解}}
\textbf{思路：}先由特征方程求出 $y(x)$，确认 $y,y'$ 在 $x\to+\infty$ 时都趋于零；也可以直接在原方程两边从 $0$ 到 $+\infty$ 积分，用初值条件求 $I$。

特征方程为
\[
r^2+2r+5=0,
\]
解得
\[
r=-1\pm 2i.
\]
故
\[
y=e^{-x}\left(C_1\cos 2x+C_2\sin 2x\right).
\]
由 $y(0)=1$ 得 $C_1=1$。求导：
\[
y'=e^{-x}\left[(-C_1+2C_2)\cos 2x+(-C_2-2C_1)\sin 2x\right].
\]
由 $y'(0)=0$ 得
\[
-C_1+2C_2=0
\quad\Longrightarrow\quad
C_2=\frac12.
\]
因此
\[
y=e^{-x}\left(\cos 2x+\frac12\sin 2x\right).
\]
由于特征根的实部为 $-1<0$，当 $x\to+\infty$ 时 $y\to0$ 且 $y'\to0$，积分收敛。

在方程 $y''+2y'+5y=0$ 两边从 $0$ 到 $+\infty$ 积分：
\[
0=\int_0^{+\infty}y''\,\mathrm{d}x+2\int_0^{+\infty}y'\,\mathrm{d}x+5\int_0^{+\infty}y\,\mathrm{d}x.
\]
于是
\[
0=\bigl[y'\bigr]_0^{+\infty}+2\bigl[y\bigr]_0^{+\infty}+5I
=-y'(0)-2y(0)+5I.
\]
代入 $y(0)=1,\ y'(0)=0$：
\[
0=-0-2+5I
\quad\Longrightarrow\quad
I=\frac25.
\]
也可直接积分验证：
\[
\int_0^{+\infty}e^{-x}\cos 2x\,\mathrm{d}x=\frac15,\qquad
\int_0^{+\infty}e^{-x}\sin 2x\,\mathrm{d}x=\frac25,
\]
故
\[
I=\frac15+\frac12\cdot\frac25=\frac25.
\]

\textbf{答案：}\boxed{I=\frac25.}

\else \answerarea[10] \fi%
